\section{Methodology}

\subsection{Research Design}

This study adopts a small-sample mixed methods design, combining descriptive results from an online questionnaire with semi-structured interview materials to examine how Asian female college students understand the relationship among campus gym space, body image pressure, and resistance training avoidance.
A mixed-methods design was chosen because the research questions require both an overview of group-level patterns and an in-depth understanding of individual experience mechanisms.
Questionnaire data can reveal which spatial perceptions and avoidance reasons are most prominent, but the aggregate export format limits individual-level inference; conversely, interviews can capture contextual nuance, yet a sample of three cannot estimate population tendencies.
By integrating both, the study addresses the gap identified in the literature review, where spatial pressure and digital appearance pressure have typically been examined separately rather than within a single analytical framework.
The original design followed a sequential explanatory mixed methods approach: the questionnaire would first establish overall trends, followed by qualitative interviews to explain those trends.
After actual data collection, the quantitative component comprised item-level aggregate export tables from the survey platform and a subsequently supplemented participant-level response matrix; the qualitative component consisted of one-on-one semi-structured interview transcripts from 3 interviewees.
Therefore, beyond descriptive statistics, this study supplements the analysis with small-sample exploratory individual-level reliability, correlation, and group comparisons, but does not conduct multivariate regression, mediation modeling, or causal inference.

This adjustment does not alter the research questions themselves, but rather the strength and mode of expression of the evidence.
Quantitative results are used to illustrate which spatial perceptions, avoidance reasons, and intervention preferences are more prominent in the sample and to preliminarily examine bivariate relationships among spatial pressure, media appearance internalization, and free-weight area avoidance; qualitative materials are used to explain the contextual mechanisms behind these numbers.
This study continues to use SEAM's approach to measuring social exercise anxiety and avoidance \parencite{levinson2013seam}, SATAQ-4's approach to measuring sociocultural appearance internalization \parencite{schaefer2015sataq4}, and research on gendered gym space \parencite{turnock2021gyms,rendall2024gym} as theoretical references, but does not claim to have used the complete original scales.

\subsection{Participants and Data Sources}

The questionnaire was distributed online via the Wenjuanxing platform, with links shared to WeChat groups and Moments of the HKUST-GZ Class of 2024, recruiting participants through convenience sampling.
Data collection took place from April 30 to May 4, 2026.
A total of 34 individuals initiated the questionnaire.
Inclusion criteria were: (a) current university student status (Question 1, 33 passed); (b) self-identified as female (Question 2, 32 passed).
The core analytical sample was therefore $N = 32$ female university students.
The age distribution was concentrated in the 18--20 range (87.50\%), with 3 individuals aged 21--23 (9.38\%) and 1 individual aged 24 and above (3.13\%).
Because the questionnaire did not separately measure ethnicity, nationality, or ethnic identity, ``Asian female college students'' in this paper primarily refers to female students enrolled at the university within the Chinese-language and Asian campus context of this study; the findings should not be generalized as the universal experience of all Asian female students.

Qualitative materials were drawn from semi-structured interviews with 3 interviewees; interview topics were consistent with the questionnaire, focusing on campus gym experiences, perceptions of the free-weight area, experiences of being watched or evaluated, social media appearance content, exercise goals, and suggestions for university support measures.
The 3 interviewees had varied experiences: one used the gym and strength training area more frequently, one had begun strength training that semester due to coursework, and one trained primarily in the dormitory and had only briefly entered the gym.
This variation facilitated comparison of how ``relatively familiar users,'' ``recent entrants,'' and ``marginal/avoiders'' differently understood the same space.

\subsection{Survey Instrument}

The questionnaire comprised 29 items.
Items 1--2 were screening questions; Items 3--9 collected age, exercise frequency, campus gym use, commonly used areas, strength training experience, and exercise goals.
Items 10--15 used 1--5 slider scales to measure spatial pressure and training anxiety in the free-weight area, including difficulty accessing the layout, discomfort when more men are present, feelings of evaluation from mirrors and public visibility, concern about incorrect form, perceptions of a male-dominated atmosphere, and reluctance to stay due to crowding and equipment concentration.
Items 16--18 measured the frequency of active avoidance of the free-weight area, reasons for avoidance, and coping strategies.
Items 19--23 measured social media use and the influence of appearance content, with Items 21--23 using 1--5 slider scales.
Item 24 measured training self-efficacy for exercise goals; Items 25--27 measured acceptance of three types of potential interventions: reducing stereotypes, increasing semi-private partitioning or visual screening, and offering beginner-friendly women's strength training workshops.

Compared with the originally planned 60--100 individual-level questionnaires and 5--6 person qualitative group interviews from the proposal, the final available sample size was smaller.
However, the subsequently supplemented text-based Excel file provided each participant's item-level responses from Item 3 onward, so this study reports both the valid sample size per item, option frequencies, slider item means, score range proportions, and individual-level construct means.
Sliders were exported by the survey platform into five segments: 1--1.8, 1.9--2.6, 2.7--3.4, 3.5--4.2, 4.3--5.
This paper merges the two upper segments (3.5--5) as the ``upper score range'' to describe the proportion of respondents who strongly endorsed a given judgment.
Individual-level construct scores were calculated as the mean of relevant items: Items 10--15 for spatial pressure/training anxiety, Items 21--23 for social media appearance internalization, Item 24 as a single-item measure of training self-efficacy, and Items 25--27 for the intervention preference index.
Because Items 25--27 each target different intervention modalities, this paper treats them as a preference item set rather than a strict reflective psychological scale.

\subsection{Qualitative Procedure and Analysis}

All 3 interviews were conducted face-to-face offline, each lasting approximately 20 minutes; interviews were audio-recorded throughout, automatically transcribed using OpenAI Whisper, and the researcher subsequently proofread and anonymized the transcripts.
Interviews followed a semi-structured protocol but allowed interviewees to elaborate based on personal experience.
Questions entered through specific scenarios, such as the feeling of first entering the gym, differences between the cardio area and the strength area, concerns about form being evaluated, whether social media influenced exercise choices, and what kind of courses or spatial support they wished the university would provide.
Interview materials were organized into a Q\&A format and anonymized; during analysis, interviewees are referred to only as P1, P2, and P3.

Qualitative analysis followed a thematic analysis approach.
The first step involved repeated reading of the 3 transcripts, marking meaning units relevant to the research questions; the second step involved initial coding around spatial visibility, beginner barrier, male-dominated atmosphere, strategic negotiation, appearance influence, and intervention needs; the third step involved merging related codes into themes and using representative quotes to test whether the themes could explain the quantitative findings.
Because the sample consisted of only 3 interviewees, the thematic analysis did not aim for saturation or broad generalization, but rather served to explain the higher levels of spatial pressure, avoidance, and intervention preferences observed in the questionnaire.

To enhance the trustworthiness of the qualitative analysis, the researcher proofread each Whisper-generated transcript against the original audio for accuracy.
The three-step thematic procedure involved iterative reading, coding, and revision rather than a single pass.
The research team included members who share the demographic background of the study population (Chinese female university students with campus gym experience), which facilitated cultural understanding of interview narratives but also required conscious reflexivity about potential confirmation bias.
Due to the time constraints of the course project, member checking was not conducted; this limitation is acknowledged.

\subsection{Data Analysis}

Quantitative data were organized and summarized using Python.
The analysis script first parsed the aggregate Excel table exported by the survey platform, breaking down items, option frequencies, slider means, and segment frequencies into structured tables; it then read the participant-level response matrix to generate a de-identified participant score table.
At the individual level, this paper calculated Cronbach's $\alpha$ for spatial pressure/training anxiety and social media appearance internalization, used Spearman's $\rho$ to test exploratory correlations between construct scores and free-weight area avoidance frequency, and used Mann--Whitney $U$ tests to compare differences between frequent and infrequent free-weight area users and between those with and without systematic strength training experience.
All $p$ values are interpreted as small-sample exploratory evidence and are not used for strong causal claims.
Qualitative data were similarly processed using Python to parse Markdown interview texts, generating Q\&A tables, thematic coding tables, and representative quote tables.

Integration of quantitative and qualitative data occurred in the results and discussion sections.
This paper first reports high-frequency phenomena from the questionnaire, then uses interviews to explain how these phenomena are experienced by interviewees as mirror visibility, male numerical dominance, fear of incorrect form, equipment knowledge barriers, and strategic entry.
This integration approach suits the current data format: bivariate statistics can indicate whether variables co-vary but cannot establish statistical causation; interview materials help form a contextualized evidence chain regarding gendered entry barriers in campus gyms.

\subsection{Ethical Considerations}

This study discusses body image, feelings of being watched, and gym discomfort, which may touch upon shame, anxiety, or self-evaluation pressure.
Both the questionnaire and interviews were based on the principles of voluntary participation, anonymous processing, and the right to withdraw at any time; participants were informed of the study's purpose before responding, and interview recordings were conducted with the interviewees' informed consent.
Interview materials are presented in this report using only de-identified quotes (P1, P2, P3), avoiding the disclosure of details that could identify individuals.
This project is a course research project and has not been reviewed by a formal institutional ethics review committee, but the research process followed the ethical guidelines for coursework; the results are used solely for course paper writing and are not intended for individual diagnosis or institutional administrative evaluation.
